% Standalone: GRASP-RAG PGR Context Formats — all four motifs
% Compile: pdflatex fig_pgr_all.tex
\documentclass{article}
\usepackage[paperwidth=24cm, paperheight=34cm,
            top=0.7cm, bottom=0.7cm, left=0.7cm, right=0.7cm]{geometry}
\usepackage{tcolorbox}
\tcbuselibrary{skins,breakable}
\usepackage{xcolor}
\usepackage[T1]{fontenc}
\usepackage{microtype}
\usepackage{tikz}
\usetikzlibrary{arrows.meta,positioning,shapes.misc}
\usepackage{array}
\pagestyle{empty}
\setlength{\parindent}{0pt}

% ── Palette ──────────────────────────────────────────────────────────────────
\definecolor{graspblue}{RGB}{37,99,235}
\definecolor{stepbg}{RGB}{239,246,255}
\definecolor{stepborder}{RGB}{147,197,253}
\definecolor{connbg}{RGB}{255,247,237}
\definecolor{connborder}{RGB}{253,186,116}
\definecolor{profilbg}{RGB}{240,253,244}
\definecolor{profilborder}{RGB}{134,239,172}
\definecolor{multibg}{RGB}{250,245,255}
\definecolor{multiborder}{RGB}{196,181,253}
\definecolor{broadbg}{RGB}{249,250,251}
\definecolor{broadborder}{RGB}{209,213,219}
\definecolor{supportbg}{RGB}{254,252,232}
\definecolor{supportborder}{RGB}{253,224,71}
\definecolor{panelbg}{RGB}{248,250,252}
\definecolor{querybg}{RGB}{30,58,138}

% ── Box types ────────────────────────────────────────────────────────────────
% Entity step (blue)
\newtcolorbox{entitybox}[1]{
  colback=stepbg, colframe=stepborder, boxrule=0.7pt, arc=3pt,
  top=2pt, bottom=2pt, left=5pt, right=4pt,
  fonttitle=\scriptsize\bfseries\sffamily\color{graspblue},
  title={\strut #1}, titlerule=0pt,
  attach boxed title to top left={yshift=-1.2mm, xshift=4mm},
  boxed title style={colback=white, colframe=stepborder, boxrule=0.5pt,
                     arc=2pt, left=2pt, right=2pt, top=0pt, bottom=0pt},
}
% Connection step (orange)
\newtcolorbox{connbox}[1]{
  colback=connbg, colframe=connborder, boxrule=0.7pt, arc=3pt,
  top=2pt, bottom=2pt, left=5pt, right=4pt,
  fonttitle=\scriptsize\bfseries\sffamily\color{orange!70!black},
  title={\strut #1}, titlerule=0pt,
  attach boxed title to top left={yshift=-1.2mm, xshift=4mm},
  boxed title style={colback=white, colframe=connborder, boxrule=0.5pt,
                     arc=2pt, left=2pt, right=2pt, top=0pt, bottom=0pt},
}
% Entity profile (green)
\newtcolorbox{profilebox}[1]{
  colback=profilbg, colframe=profilborder, boxrule=0.7pt, arc=3pt,
  top=2pt, bottom=2pt, left=5pt, right=4pt,
  fonttitle=\scriptsize\bfseries\sffamily\color{green!50!black},
  title={\strut #1}, titlerule=0pt,
  attach boxed title to top left={yshift=-1.2mm, xshift=4mm},
  boxed title style={colback=white, colframe=profilborder, boxrule=0.5pt,
                     arc=2pt, left=2pt, right=2pt, top=0pt, bottom=0pt},
}
% Multi-entity block (purple)
\newtcolorbox{multibox}[1]{
  colback=multibg, colframe=multiborder, boxrule=0.7pt, arc=3pt,
  top=2pt, bottom=2pt, left=5pt, right=4pt,
  fonttitle=\scriptsize\bfseries\sffamily\color{violet!70!black},
  title={\strut #1}, titlerule=0pt,
  attach boxed title to top left={yshift=-1.2mm, xshift=4mm},
  boxed title style={colback=white, colframe=multiborder, boxrule=0.5pt,
                     arc=2pt, left=2pt, right=2pt, top=0pt, bottom=0pt},
}
% Supporting context (yellow)
\newtcolorbox{supportbox}{
  colback=supportbg, colframe=supportborder, boxrule=0.6pt, arc=3pt,
  top=2pt, bottom=2pt, left=5pt, right=4pt,
}
% Broad/flat (gray)
\newtcolorbox{broadbox}[1]{
  colback=broadbg, colframe=broadborder, boxrule=0.7pt, arc=3pt,
  top=2pt, bottom=2pt, left=5pt, right=4pt,
  fonttitle=\scriptsize\bfseries\sffamily\color{gray!70},
  title={\strut #1}, titlerule=0pt,
  attach boxed title to top left={yshift=-1.2mm, xshift=4mm},
  boxed title style={colback=white, colframe=broadborder, boxrule=0.5pt,
                     arc=2pt, left=2pt, right=2pt, top=0pt, bottom=0pt},
}

% ── Shorthand ────────────────────────────────────────────────────────────────
\newcommand{\codekw}[1]{\texttt{\color{graspblue}#1}}
\newcommand{\fn}[1]{\texttt{\footnotesize\color{gray!70}#1}}
\newcommand{\passagemark}{\textit{\color{gray}\scriptsize [source passage\ldots]}}
\newcommand{\paneltitle}[3]{%
  \begin{tcolorbox}[
    colback=querybg, colframe=querybg,
    boxrule=0pt, arc=4pt, top=4pt, bottom=4pt, left=6pt, right=6pt]
    {\small\bfseries\sffamily\color{white} #1}
    \hfill{\tiny\sffamily\color{white!70}\texttt{#2}}\\[1pt]
    {\scriptsize\sffamily\color{white!85}\textit{``#3''}}
  \end{tcolorbox}%
}

\begin{document}

% ── Page title ────────────────────────────────────────────────────────────────
{\centering
{\large\bfseries\sffamily GRASP-RAG: Path-Guided Reasoning (PGR) — Context Formats}\\[2pt]
{\small\sffamily\color{gray}
  One shared generation prompt; motif-specific structure lives in the \textbf{context}.
  Context builders: \fn{\_bridge\_context}, \fn{\_single\_entity\_context},
  \fn{\_multi\_entity\_context}; broad falls through to LightRAG hybrid.}
\par}

\vspace{6pt}
\noindent\rule{\linewidth}{0.5pt}
\vspace{4pt}

% ══════════════════════════════════════════════════════════════════════════════
%  TWO COLUMN LAYOUT
% ══════════════════════════════════════════════════════════════════════════════
\begin{minipage}[t]{0.487\textwidth}

%% ─────────────────────────────────────────────────────────────────────────────
%%  (a) BRIDGE — _bridge_context
%% ─────────────────────────────────────────────────────────────────────────────
\paneltitle{(a) Bridge Query}{\_bridge\_context}{What connects Rational Grazing to Polyface Farm?}

\vspace{4pt}

{\scriptsize\sffamily\color{gray!60}
\textit{Path selected by edge-keyword relevance to query (not shortest-path only).
Remaining chunk budget filled via graph-guided re-query using best path's keywords.}}

\vspace{4pt}

\begin{entitybox}{Step 1 — Entity: Rational Grazing \normalfont\color{gray}(method)}
\scriptsize\sffamily
\textbf{Description:} Systematic approach to managing pasture through controlled
grazing periods; maximises productivity and recovery\ldots\\
\passagemark
\end{entitybox}

\vspace{3pt}

\begin{connbox}{Step 2 — Connection: Rational Grazing $\boldsymbol{\rightarrow}$ \textbf{Cows}}
\scriptsize\sffamily
\textbf{Keywords:} livestock management, rotational grazing, animal health\\
\passagemark
\end{connbox}

\vspace{3pt}

\begin{entitybox}{Step 3 — Entity: \textbf{Cows} \normalfont\color{gray}(livestock)}
\scriptsize\sffamily
\textbf{Description:} Central to regenerative systems; Polyface Farm uses cows
as primary animal in rotational schedule driven by rational grazing logic\ldots\\
\passagemark
\end{entitybox}

\vspace{3pt}

\begin{connbox}{Step 4 — Connection: \textbf{Cows} $\boldsymbol{\rightarrow}$ Polyface Farm}
\scriptsize\sffamily
\textbf{Keywords:} rotational livestock, farm operations, pasture management\\
\passagemark
\end{connbox}

\vspace{3pt}

\begin{entitybox}{Step 5 — Entity: Polyface Farm \normalfont\color{gray}(organisation)}
\scriptsize\sffamily
\textbf{Description:} Joel Salatin's regenerative farm in Virginia; multi-species
rotational model with cows, pigs, and poultry\ldots\\
\passagemark
\end{entitybox}

\vspace{3pt}

\begin{supportbox}
\scriptsize\sffamily\color{gray!80}
\textbf{[Supporting Context]} \textit{(alternative paths + graph-guided re-query chunks)}\\
\ldots additional evidence for the specific connection\ldots
\end{supportbox}

\vspace{4pt}
{\scriptsize\sffamily\color{graspblue!80}
\textbf{Effect:} Bridge entity (\textbf{Cows}) is structurally unavoidable at Step 3;
LLM must traverse the path to answer.}

\vspace{10pt}

%% ─────────────────────────────────────────────────────────────────────────────
%%  (b) SINGLE-ENTITY — _single_entity_context
%% ─────────────────────────────────────────────────────────────────────────────
\paneltitle{(b) Single-Entity Query}{\_single\_entity\_context}{What is Polyface Farm?}

\vspace{4pt}

{\scriptsize\sffamily\color{gray!60}
\textit{Anchor = first node of first returned path. Direct connections drawn from
all returned paths (up to 10). Secondary chunks filled by VDB re-query.}}

\vspace{4pt}

\begin{profilebox}{Entity Profile: Polyface Farm}
\scriptsize\sffamily
\textbf{Type:} organisation\\
\textbf{Description:} Polyface Farm is Joel Salatin's regenerative farm in
Swoope, Virginia. Known for multi-species rotational grazing\ldots\\[3pt]
\textbf{Connected to:}\\
\quad $\bullet$ \textbf{Joel Salatin}: ownership, farm management\\
\quad $\bullet$ \textbf{Cows}: rotational livestock, primary grazer\\
\quad $\bullet$ \textbf{Rational Grazing}: grazing method, pasture recovery\\
\quad $\bullet$ \textbf{Chickens}: multi-species rotation, following cows\\
\quad $\bullet$ \textbf{Holistic Management}: land stewardship, philosophy\\[3pt]
\textbf{Source Passages:}\\
\passagemark\quad \passagemark\quad \passagemark
\end{profilebox}

\vspace{4pt}
{\scriptsize\sffamily\color{green!50!black!90}
\textbf{Effect:} LLM sees a rich entity profile plus all direct relations — optimised
for definitional / factual single-entity questions.}

\end{minipage}
\hfill
\begin{minipage}[t]{0.487\textwidth}

%% ─────────────────────────────────────────────────────────────────────────────
%%  (c) MULTI-ENTITY — _multi_entity_context
%% ─────────────────────────────────────────────────────────────────────────────
\paneltitle{(c) Multi-Entity Query}{\_multi\_entity\_context}{Compare Rational Grazing, Holistic Management, and Polyface Farm.}

\vspace{4pt}

{\scriptsize\sffamily\color{gray!60}
\textit{Unique nodes from all Steiner-tree paths collected as anchors.
Per-entity profiles built first; connection map built from all path edges.}}

\vspace{4pt}

\begin{multibox}{Entity 1: Rational Grazing \normalfont\color{gray}(method)}
\scriptsize\sffamily
\textbf{Description:} Systematic pasture management through controlled grazing
and recovery periods; maximises productivity\ldots\\
\passagemark
\end{multibox}

\vspace{3pt}

\begin{multibox}{Entity 2: Holistic Management \normalfont\color{gray}(framework)}
\scriptsize\sffamily
\textbf{Description:} Decision-making framework developed by Allan Savory;
integrates social, economic, and ecological factors\ldots\\
\passagemark
\end{multibox}

\vspace{3pt}

\begin{multibox}{Entity 3: Polyface Farm \normalfont\color{gray}(organisation)}
\scriptsize\sffamily
\textbf{Description:} Joel Salatin's flagship regenerative farm; demonstrates
multi-species rotation and integrated land management\ldots\\
\passagemark
\end{multibox}

\vspace{3pt}

\begin{multibox}{Steiner intermediate: Cows \normalfont\color{gray}(livestock)}
\scriptsize\sffamily
\textbf{Description:} Primary grazing animal linking all three above entities\ldots\\
\passagemark
\end{multibox}

\vspace{3pt}

\begin{tcolorbox}[colback=multibg!60, colframe=multiborder, boxrule=0.7pt,
  arc=3pt, top=2pt, bottom=2pt, left=5pt, right=4pt]
\scriptsize\sffamily\textbf{\color{violet!70!black}[Connections]}\\
\quad $\bullet$ Rational Grazing $\leftrightarrow$ Cows: \textit{livestock management, rotation}\\
\quad $\bullet$ Holistic Management $\leftrightarrow$ Cows: \textit{herd impact, land regeneration}\\
\quad $\bullet$ Cows $\leftrightarrow$ Polyface Farm: \textit{rotational livestock, operations}\\
\quad $\bullet$ Rational Grazing $\leftrightarrow$ Holistic Management: \textit{land philosophy}
\end{tcolorbox}

\vspace{4pt}
{\scriptsize\sffamily\color{violet!70!black!90}
\textbf{Effect:} LLM sees per-entity profiles and an explicit connection map — optimised
for comparison and multi-hop integrative questions.}

\vspace{10pt}

%% ─────────────────────────────────────────────────────────────────────────────
%%  (d) BROAD — LightRAG hybrid fallback (no custom context)
%% ─────────────────────────────────────────────────────────────────────────────
\paneltitle{(d) Broad Query}{LightRAG hybrid fallback}{What are the key debates in regenerative agriculture?}

\vspace{4pt}

{\scriptsize\sffamily\color{gray!60}
\textit{No custom PGR format. Code calls \fn{rag.aquery(query, param=SearchParam(mode="hybrid"))}.
Context is built by LightRAG: VDB chunk retrieval + keyword graph traversal, flat.}}

\vspace{4pt}

\begin{broadbox}{VDB chunks (vector similarity)}
\scriptsize\sffamily\color{gray!80}
Passage 1: Regenerative agriculture encompasses practices like cover cropping,
reduced tillage, and rotational grazing\ldots\\[2pt]
Passage 2: Critics argue that regenerative claims lack standardised
certification and measurable carbon sequestration benchmarks\ldots\\[2pt]
Passage 3: Holistic Management's savanna hypothesis has been contested by
ecologists who dispute the scalability of herd impact\ldots
\end{broadbox}

\vspace{3pt}

\begin{broadbox}{Graph keywords (LightRAG high-level)}
\scriptsize\sffamily\color{gray!80}
\textit{keywords: soil health, carbon sequestration, cover crops,
rotational grazing, biodiversity, no-till farming, certification\ldots}
\end{broadbox}

\vspace{3pt}

{\scriptsize\sffamily\color{gray!60}
\textit{No step-by-step labels. No entity profiles. No connection map.
Flat, unordered retrieval — same as LightRAG's own output.}}

\vspace{4pt}
\begin{tcolorbox}[colback=gray!8, colframe=gray!30, boxrule=0.6pt, arc=3pt,
  top=3pt, bottom=3pt, left=5pt, right=4pt]
\scriptsize\sffamily\color{gray!80}
\textbf{Why:} Broad queries have no specific entity anchors, so graph
traversal motifs don't apply. GRASP-RAG delegates to LightRAG hybrid,
which is optimised for diffuse thematic retrieval.
Win rate vs LightRAG on broad queries is 69.6\% — not structural.
\end{tcolorbox}

\end{minipage}

\vspace{6pt}
\noindent\rule{\linewidth}{0.5pt}
\vspace{3pt}

% ── Bottom legend: colour key + prompt note ───────────────────────────────────
\begin{minipage}[t]{0.52\textwidth}
{\scriptsize\sffamily
\textbf{Colour key:}
\begin{tabular}{@{}l@{\;=\;}l@{\quad}l@{\;=\;}l@{}}
  \textcolor{graspblue}{\rule{6pt}{6pt}} & Entity step (blue) &
  \textcolor{orange!80!black}{\rule{6pt}{6pt}} & Connection step (orange) \\
  \textcolor{green!50!black}{\rule{6pt}{6pt}} & Entity profile (green) &
  \textcolor{violet!70!black}{\rule{6pt}{6pt}} & Multi-entity / connections (purple) \\
  \textcolor{yellow!70!black}{\rule{6pt}{6pt}} & Supporting context (yellow) &
  \textcolor{gray!60}{\rule{6pt}{6pt}} & Broad / flat (gray) \\
\end{tabular}}
\end{minipage}
\hfill
\begin{minipage}[t]{0.44\textwidth}
\begin{tcolorbox}[colback=graspblue!5, colframe=graspblue!30, boxrule=0.6pt,
  arc=3pt, top=3pt, bottom=3pt, left=5pt, right=4pt]
\scriptsize\sffamily\color{graspblue!80}
\textbf{One shared prompt for all motifs} (\fn{\_RESPONSE\_PROMPT}).
Structure is entirely in the context, not in the prompt instruction.
Each motif's context builder imposes the topology; the LLM reads it as-is.
\end{tcolorbox}
\end{minipage}

\end{document}
